\documentclass[ordinary]{BMSTU-IU8}

\KOMAoptions{fontsize=14pt}

\usepackage{fontspec}
\setmainfont{Times New Roman}

\geometry{a4paper,left=30mm,right=15mm,top=20mm,bottom=20mm,footskip=10mm}
\onehalfspacing
\setlength{\parskip}{0pt}
\setlength{\parindent}{1.25cm}
\raggedbottom
\clubpenalty=10000
\widowpenalty=10000

\newlength{\rpzheadinggap}
\setlength{\rpzheadinggap}{10mm}

\addbibresource{references.bib}
\DefineBibliographyStrings{russian}{urlseen={дата обращения:}}
\defbibenvironment{bibliography}
  {\list{}
    {\setlength{\leftmargin}{0pt}
     \setlength{\itemindent}{\parindent}
     \setlength{\itemsep}{\bibitemsep}
     \setlength{\parsep}{\bibparsep}}}
  {\endlist}
  {\item\printtext[labelnumberwidth]
    {\printfield{labelprefix}\printfield{labelnumber}}\addspace}

\titleformat{\structure}[block]
  {\normalfont\bfseries\fontsize{16pt}{24pt}\selectfont\centering
   \hyphenpenalty=10000\exhyphenpenalty=10000}{}{0pt}{}
\titlespacing{\structure}{0pt}{0pt}{\rpzheadinggap}
\titleformat{\section}[block]
  {\normalfont\bfseries\fontsize{16pt}{24pt}\selectfont
   \hyphenpenalty=10000\exhyphenpenalty=10000}{\thesection}{1em}{}
\titlespacing{\section}{1.25cm}{0pt}{\rpzheadinggap}
\titleformat{\subsection}[block]
  {\normalfont\bfseries\hyphenpenalty=10000\exhyphenpenalty=10000}
  {\thesubsection}{1em}{}
\titlespacing{\subsection}{1.25cm}{\rpzheadinggap}{\rpzheadinggap}

\begin{document}

\section{Аналитическая часть}

В разделе описываются состав планетной сцены и её параметры. Приводится математическое описание объектов
и движения планет. Рассматриваются методы построения и отображения
трёхмерной сцены, проводится их сравнение и обосновывается выбор для
разрабатываемого конструктора. По результатам анализа формулируется
постановка задачи и уточняются ограничения модели.

\subsection{Анализ предметной области}
\label{sec:domain}

В работе рассматривается задача трёхмерной визуализации планетных систем. Моделируемая система включает центральную звезду и планеты, движущиеся вокруг неё по орбитам. Примером такой структуры служит Солнечная система,
планеты которой обращаются вокруг Солнца~\cite{nasa-solar-system}.
Разрабатываемый конструктор предназначен для создания визуализации собственной
планетной системы и изменения её параметров через пользовательский интерфейс.

Для формирования изображения трёхмерной сцены необходимо описать геометрию
объектов, свойства их поверхностей, источники света и камеру. Геометрия
определяет форму и положение объектов, материал задаёт особенности
взаимодействия поверхности со светом, а камера определяет условия
наблюдения сцены~\cite[разд.~1.2]{pharr2023}. 

В состав моделируемой сцены входят:
\begin{itemize}
    \item центральная звезда --- объект, присутствующий на сцене при запуске
    программы и служащий источником света с настраиваемым цветом и
    параметрами освещения;
    \item планеты --- добавляемые пользователем тела, представленные сферами
    различного радиуса и имеющие настраиваемые свойства материала;
    \item камеры --- средства наблюдения сцены с изменяемым положением.
\end{itemize}

Геометрическая модель звезды и источник освещения выполняют разные функции:
первая определяет видимое представление звезды, второй --- освещение
остальных тел. Пользователь должен иметь возможность изменять цвет
источника и параметры освещения. Внешний вид планет определяется
совместным влиянием геометрии, свойств материала и освещения; поэтому
параметры материала и источника света задаются раздельно.

При описании движения необходимо различать перемещение планеты по орбите
и вращение вокруг собственной оси. В первом случае изменяется положение
центра тела относительно звезды, во втором --- ориентация тела относительно
собственного центра. В конструкторе скорости этих движений могут
изменяться независимо. Пользователь также может изменять размеры орбит
планет. Радиус планеты характеризует размер самого тела
и задаётся отдельно от размеров его орбиты.

Пользователь должен иметь возможность наблюдать движение планет и результат
изменения параметров сцены в процессе работы программы. По этой причине формируемое
изображение должно обновляться при изменении положения и ориентации тел,
настроек материалов, освещения и камеры.

\printbibliography

\end{document}
